\section{习题}

\subsection{基础过关题}

\begin{enumerate}

\item \textbf{（单项选择题）} 某机字长 8 位，定点整数采用补码表示。下列叙述中正确的是
\begin{center}
\begin{tabular}{ll}
(A) 最小负数为 $-127$，编码为 \texttt{1000 0000} & (B) 最小负数为 $-128$，编码为 \texttt{1000 0000} \\[4pt]
(C) 编码 \texttt{1111 1111} 表示 $-0$ & (D) 最大正数为 $+128$，编码为 \texttt{0111 1111}
\end{tabular}
\end{center}

\ifshowsolutions\par\noindent\textbf{解析：} 8 位补码的表示范围为 $[-2^{7},\,2^{7}-1]=[-128,\,127]$。编码 \texttt{1000 0000} 的符号位为 1，按“取反加 1”求真值：$\sim\!\texttt{1000 0000}+1=\texttt{1000 0000}=128$，故它表示 $-128$，这正是补码比原码多出的那个最小负数。(A) 中 $-127$ 的编码应为 \texttt{1000 0001}；(C) 中 \texttt{1111 1111} 表示 $-1$，补码零唯一，不存在 $-0$；(D) 中 $+128$ 已超出 8 位补码范围，\texttt{0111 1111} 表示 $+127$。\boxed{\text{\textbf{答案：}(B)}}\else \answerarea \fi

\item \textbf{（填空题）} 8 位补码 \texttt{1110 0110} 的真值为 \underline{\qquad\qquad}；将它符号扩展为 16 位后，得到的 16 位机器数为 \underline{\qquad\qquad}（用十六进制表示）。

\ifshowsolutions\par\noindent\textbf{解析：} 符号位为 1，是负数。对 \texttt{1110 0110} 连同符号位取反再加 1：
\[
\sim\!\texttt{1110 0110}=\texttt{0001 1001},\qquad
\texttt{0001 1001}+1=\texttt{0001 1010}=26,
\]
所以真值为 $-26$。符号扩展时高位全部补原符号位 1，低 8 位保持不变：
\[
\texttt{1111 1111 1110 0110}=\texttt{FFE6H}.
\]
也可直接验证：$\texttt{FFE6H}$ 取反加 1 得 $\texttt{001A}\text{H}=26$，确为 $-26$。\boxed{\text{\textbf{答案：}$-26$；\texttt{FFE6H}}}\else \answerarea \fi

\item \textbf{（单项选择题）} 8 位补码 $[x]_{\text{补}}=\texttt{7BH}$，$[y]_{\text{补}}=\texttt{0AH}$。执行 $x+y$ 后，结果的机器数以及进位标志 CF、溢出标志 OF 分别为
\begin{center}
\begin{tabular}{ll}
(A) \texttt{85H}，CF$=0$，OF$=1$ & (B) \texttt{85H}，CF$=1$，OF$=0$ \\[4pt]
(C) \texttt{05H}，CF$=0$，OF$=1$ & (D) \texttt{85H}，CF$=1$，OF$=1$
\end{tabular}
\end{center}

\ifshowsolutions\par\noindent\textbf{解析：} 先求真值：$\texttt{7BH}=0111\,1011_2=+123$，$\texttt{0AH}=0000\,1010_2=+10$，真值和为 $133$，已超过 8 位补码最大值 $127$，故必然溢出。按位相加：
\[
\begin{array}{r}
0111\,1011\\
+\,0000\,1010\\ \hline
1000\,0101
\end{array}
\]
结果为 \texttt{85H}。符号位（第 7 位）向外的进位 $C_7=0$，最高数值位（第 6 位）向符号位的进位 $C_6=1$，故
\[
\text{OF}=C_7\oplus C_6=0\oplus1=1.
\]
CF 取符号位的进位输出，即 CF$=C_7=0$。两个正数相加得到符号位为 1 的结果，也直接说明发生了正溢出。\boxed{\text{\textbf{答案：}(A)}}\else \answerarea \fi

\item \textbf{（单项选择题）} 下列关于 IEEE 754 单精度浮点数特殊编码的叙述中，正确的是
\begin{center}
\begin{tabular}{ll}
(A) \texttt{7F80 0000H} 表示最大规格化正数 & (B) \texttt{FF80 0000H} 表示 $-\infty$ \\[4pt]
(C) \texttt{0000 0001H} 表示最小规格化正数 & (D) \texttt{7FC0 0000H} 表示 $+0$
\end{tabular}
\end{center}

\ifshowsolutions\par\noindent\textbf{解析：} 单精度共 32 位：$S$ 占 1 位，阶码 $E$ 占 8 位，尾数域 $M$ 占 23 位。(A) 中阶码为 \texttt{1111 1111}、尾数域全 0，表示 $+\infty$，不是最大规格化数；(B) 中 $S=1$、阶码全 1、尾数域全 0，表示 $-\infty$，正确；(C) 中阶码全 0、尾数域非 0，是非规格化数，其值为 $2^{-149}$，最小规格化正数应为 $1.0\times2^{-126}$（编码 \texttt{0080 0000H}）；(D) 中阶码全 1、尾数域非 0，表示 NaN，$+0$ 的编码是全 0。\boxed{\text{\textbf{答案：}(B)}}\else \answerarea \fi

\item \textbf{（填空题）} 十进制数 $-0.15625$ 的 IEEE 754 单精度浮点数编码为 \underline{\qquad\qquad}（用十六进制表示）。

\ifshowsolutions\par\noindent\textbf{解析：} 先化为二进制：
\[
0.15625=\frac{5}{32}=0.00101_2=1.01_2\times2^{-3}.
\]
于是 $S=1$，真指数 $e=-3$，阶码
\[
E=e+127=-3+127=124=\texttt{0111 1100}_2.
\]
尾数域取 $1.01_2$ 去掉隐藏位 1 后的小数部分 \texttt{01}，右侧补 0 至 23 位：
\[
M=\texttt{0100 0000 0000 0000 0000 000}.
\]
拼接得
\[
\texttt{1\ 0111 1100\ 0100 0000 0000 0000 0000 0000}
=\texttt{BE20 0000H}.
\]
反查：$S=1$，$E=124$，$e=-3$，尾数真值 $1.01_2=1.25$，故值为 $-1.25\times2^{-3}=-0.15625$。\boxed{\text{\textbf{答案：}\texttt{BE20 0000H}}}\else \answerarea \fi

\item \textbf{（简答题）} 某 $n$ 位 ALU 用补码加法器实现减法 $A-B$：把 $B$ 的各位取反后送入加法器，并令最低位进位输入 $C_0=1$。说明这样做的原理，并指出此时加法器最高位进位输出 $C_n$ 与无符号减法借位标志的关系。

\ifshowsolutions\par\noindent\textbf{解析：} 在 $n$ 位模 $2^n$ 意义下，$B$ 的补码（即 $-B$ 的机器数）为 $\overline{B}+1$，其中 $\overline{B}$ 表示逐位取反。因此
\[
A-B\equiv A+\overline{B}+1\pmod{2^n}.
\]
把 $\overline{B}$ 送入加法器、再令 $C_0=1$，恰好补上了取反后缺少的那个 1，加法器输出的 $n$ 位结果就是 $A-B$ 的补码。减法无需另建减法电路，只需在 B 端加一级反相器并把 $C_0$ 置 1。

对无符号数，若 $A\geqslant B$，则 $A+\overline{B}+1\geqslant 2^n$，最高位产生进位，$C_n=1$；若 $A<B$，则结果不足 $2^n$，$C_n=0$，说明需要向高位借位。因此无符号减法的借位标志为
\[
\text{borrow}=\overline{C_n}.
\]
有符号溢出仍由 $\text{OF}=C_n\oplus C_{n-1}$ 判断；它与借位是两件不同的事，不能混用。\boxed{\text{\textbf{答案：}$A-B\equiv A+\overline{B}+1$，$C_0=1$；借位 $=\overline{C_n}$}}\else \answerarea \fi

\end{enumerate}

\subsection{真题风格与综合训练}

\begin{enumerate}[resume]

\item \textbf{（真题风格·单项选择题）} 在 32 位机器中，某变量的机器数为 \texttt{BF80 0000H}。若分别按 \texttt{int}（32 位补码）和 \texttt{float}（IEEE 754 单精度）解释，则对应的值依次为
\begin{center}
\begin{tabular}{ll}
(A) $-1082130432$，$-1.0$ & (B) $-1082130432$，$-1.5$ \\[4pt]
(C) $-1073741824$，$-1.0$ & (D) $-1082130432$，$-2.0$
\end{tabular}
\end{center}

\ifshowsolutions\par\noindent\textbf{解析：} 先按 \texttt{int} 解释。\texttt{BF80 0000H} 的符号位为 1，对其取反加 1：
\[
\sim\!\texttt{BF80 0000H}=\texttt{407F FFFFH},\qquad
\texttt{407F FFFFH}+1=\texttt{4080 0000H}.
\]
而
\[
\texttt{4080 0000H}=2^{30}+2^{23}=1073741824+8388608=1082130432,
\]
故 \texttt{int} 值为 $-1082130432$。

再按 \texttt{float} 解释：最高位 $S=1$；其后 8 位为 \texttt{0111 1111}$=127$，真指数 $e=127-127=0$；尾数域 23 位全 0，尾数真值为 $1.0$。所以值为
\[
(-1)^1\times1.0\times2^{0}=-1.0.
\]
可见同一串二进制位，类型不同则含义完全不同。\boxed{\text{\textbf{答案：}(A)}}\else \answerarea[8] \fi

\item \textbf{（真题风格·单项选择题）} 设单精度浮点数 $a=1.0$，$b=2^{-24}$，二者均为规格化数。按 IEEE 754 默认舍入模式（就近舍入到偶数）计算 $a+b$，所得结果的机器数是
\begin{center}
\begin{tabular}{ll}
(A) \texttt{3F80 0000H} & (B) \texttt{3F80 0001H} \\[4pt]
(C) \texttt{3F80 0002H} & (D) \texttt{4000 0000H}
\end{tabular}
\end{center}

\ifshowsolutions\par\noindent\textbf{解析：} $a=1.0=1.0_2\times2^{0}$，编码为 \texttt{3F80 0000H}；$b=2^{-24}$ 的阶码 $E=-24+127=103=\texttt{0110 0111}_2$，编码为 \texttt{3380 0000H}。

两数相加须先对阶。$a$ 的阶码为 0，$b$ 的阶码为 $-24$，将 $b$ 的尾数右移 24 位后与 $a$ 相加，精确和为
\[
1+2^{-24}.
\]
在 $[1,2)$ 这一 binade 内，相邻单精度数的间距（ulp）为 $2^{-23}$。因此精确和恰好落在 $1.0$ 与 $1+2^{-23}$ 的正中间：
\[
1+2^{-24}=\frac{1.0+(1+2^{-23})}{2}.
\]
这是“平局”。就近舍入到偶数要求选尾数最低有效位为 0 的那个，即 $1.0$（其尾数域全 0），而不是 $1+2^{-23}$（尾数域末位为 1）。故结果为 \texttt{3F80 0000H}。若采用向 $+\infty$ 舍入，才会得到 \texttt{3F80 0001H}。\boxed{\text{\textbf{答案：}(A)}}\else \answerarea[8] \fi

\item \textbf{（综合应用题）} 设单精度浮点数 $X$ 的机器数为 \texttt{C128 0000H}，$Y$ 的机器数为 \texttt{4050 0000H}。两数均为规格化数，运算采用就近舍入到偶数，并保留保护位、舍入位和粘滞位；本例不涉及上溢、下溢和特殊值。

请完成：(1) 分别写出 $X$、$Y$ 的十进制真值；(2) 按“对阶—尾数加减—规格化—舍入—溢出判断”计算 $X+Y$，写出中间尾数形式；(3) 给出最终结果的 IEEE 754 单精度编码。

\ifshowsolutions\par\noindent\textbf{解析：} \textbf{(1) 解码。}

$X=\texttt{C128 0000H}$ 的二进制为
\[
\texttt{1\ 1000 0010\ 0101 0000 0000 0000 0000 0000}.
\]
故 $S_X=1$，$E_X=\texttt{1000 0010}_2=130$，真指数 $e_X=130-127=3$，尾数真值为 $1.0101_2$。于是
\[
X=-1.0101_2\times2^{3}=-1010.1_2=-10.5.
\]

$Y=\texttt{4050 0000H}$ 的二进制为
\[
\texttt{0\ 1000 0000\ 1010 0000 0000 0000 0000 0000},
\]
故 $S_Y=0$，$E_Y=128$，$e_Y=1$，尾数真值为 $1.101_2$。于是
\[
Y=+1.101_2\times2^{1}=11.01_2=3.25.
\]

\textbf{(2) 浮点加法。}

对阶：$\Delta E=|3-1|=2$，小阶 $Y$ 向大阶 $X$ 看齐，$Y$ 的尾数算术右移 2 位，阶码变为 3：
\[
Y=0.01101_2\times2^{3}
\]
（附加位全为 0，无精度损失）。尾数相减：
\[
(-1.0101_2)+(0.01101_2)=-(1.01010_2-0.01101_2)
=-0.11101_2.
\]
即中间结果为 $-0.11101_2\times2^{3}$。

规格化：尾数绝对值小于 1，需左规 1 位，尾数变为 $-1.1101_2$，阶码减 1 得 $e=2$。

舍入：左规过程中没有丢弃非零低位，保护位、舍入位、粘滞位均为 0，结果精确，无需舍入。

溢出判断：$e=2$ 位于单精度规格化数指数范围 $[-126,\,127]$ 内，无上溢或下溢。

\textbf{(3) 编码。} 最终值为
\[
-1.1101_2\times2^{2}=-111.01_2=-7.25.
\]
此时 $S=1$，$E=2+127=129=\texttt{1000 0001}_2$，尾数域为 \texttt{1101} 后补 19 个 0：
\[
\texttt{1\ 1000 0001\ 1101 0000 0000 0000 0000 0000}
=\texttt{C0E8 0000H}.
\]
验证：$1.1101_2=1.8125$，$1.8125\times4=7.25$，符号为负，与 $-10.5+3.25=-7.25$ 一致。\boxed{\text{\textbf{答案：}$X=-10.5$，$Y=3.25$；$X+Y=-7.25$，编码 \texttt{C0E8 0000H}}}\else \answerarea[10] \fi

\item \textbf{（综合应用题）} 某 8 位 ALU 用补码完成加、减运算。操作数 $A=\texttt{70H}$，$B=\texttt{90H}$。执行减法时，ALU 将 $B$ 逐位取反后送入加法器，并令最低位进位输入 $C_0=1$。设 CF 为加法器最高位进位输出（减法时按“借位”解释，借位 $=\overline{C_7}$）；OF $=C_7\oplus C_6$，其中 $C_6$ 是最高数值位向符号位的进位；SF 取结果最高位，ZF 在结果 8 位全 0 时为 1。(1) $A+B$ 的结果及 CF、OF、SF、ZF；(2) $A-B$ 的结果及 CF（借位）、OF、SF、ZF；(3) 解释为什么第 (1) 问中 CF$=1$ 而 OF$=0$。

\ifshowsolutions\par\noindent\textbf{解析：} 先写出机器数：
\[
A=\texttt{70H}=0111\,0000_2,\qquad
B=\texttt{90H}=1001\,0000_2.
\]
按补码解释，$A=+112$，$B=-112$；按无符号数解释，$A=112$，$B=144$。

\textbf{(1) $A+B$。} 直接相加：
\[
\begin{array}{r}
0111\,0000\\
+\,1001\,0000\\ \hline
1\,0000\,0000
\end{array}
\]
丢弃模 $2^8$ 的进位后，结果为 \texttt{00H}。最高位进位 $C_7=1$，次高位进位 $C_6=1$，故
\[
\text{CF}=C_7=1,\qquad \text{OF}=C_7\oplus C_6=1\oplus1=0.
\]
结果全 0，故 ZF$=1$，SF$=0$。从有符号数看，$112+(-112)=0$，没有溢出；从无符号数看，$112+144=256$，产生进位，故 CF$=1$。

\textbf{(2) $A-B$。} 先取反：$\overline{B}=\texttt{6FH}=0110\,1111_2$。再令 $C_0=1$：
\[
A+\overline{B}+1
=0111\,0000_2+0110\,1111_2+1
=1110\,0000_2=\texttt{E0H}.
\]
逐位观察进位：第 6 位向第 7 位有进位，$C_6=1$；第 7 位向外无进位，$C_7=0$。因此
\[
\text{OF}=C_7\oplus C_6=0\oplus1=1.
\]
减法借位为 $\overline{C_7}=1$。结果最高位为 1，SF$=1$；结果非 0，ZF$=0$。从有符号数看，$112-(-112)=224>127$，发生正溢出，回绕后的 \texttt{E0H} 表示 $-32$；从无符号数看，$112-144$ 需要借位，故借位标志为 1。

\textbf{(3) 两个标志的含义不同。} CF 记录的是把操作数当作无符号数时最高位是否产生进位（减法时为是否借位），它只关心模 $2^8$ 之外的进位；OF 记录的是把操作数当作补码有符号数时结果是否超出 $[-128,\,127]$，由 $C_7\oplus C_6$ 判断。第 (1) 问的无符号和 $112+144=256$ 超出了 8 位无符号范围，故 CF$=1$；但有符号和 $112+(-112)=0$ 完全在范围内，故 OF$=0$。两者互不蕴含。

\boxed{\text{\textbf{答案：}(1) \texttt{00H}: CF=1, OF=0, SF=0, ZF=1; (2) \texttt{E0H}: CF=1, OF=1, SF=1, ZF=0}}\else \answerarea[10] \fi

\end{enumerate}
